\documentclass[9pt,a4paper]{extarticle}

\usepackage[utf8]{inputenc}
\usepackage[T5]{fontenc}
\usepackage[vietnamese]{babel}
\usepackage[a4paper,top=6mm,bottom=6mm,left=6mm,right=6mm,columnsep=4mm]{geometry}
\usepackage{multicol}
\usepackage{enumitem}
\usepackage{titlesec}
\usepackage{microtype}
\usepackage{xcolor}
\usepackage{lmodern}
\usepackage{parskip}

\pagestyle{empty}
\setlength{\parindent}{0pt}
\setlength{\parskip}{1pt}
\setlength{\columnseprule}{0pt}

\linespread{0.92}

\titleformat{\section}{\bfseries\fontsize{9.2}{10}\selectfont}{}{0pt}{}
\titleformat{\subsection}{\bfseries\fontsize{8.6}{9.2}\selectfont}{}{0pt}{}
\titlespacing*{\section}{0pt}{2pt}{1pt}
\titlespacing*{\subsection}{0pt}{1pt}{0.5pt}

\setlist[itemize]{leftmargin=*,topsep=1pt,itemsep=0.5pt,parsep=0pt,partopsep=0pt}
\setlist[enumerate]{leftmargin=*,topsep=1pt,itemsep=0.5pt,parsep=0pt,partopsep=0pt}

\begin{document}
\small
\begin{center}
{\bfseries\fontsize{10}{11}\selectfont Sắp xếp theo độ liên quan: Lý thuyết \& Ứng dụng}\\[-1pt]
\end{center}

\vspace{-3pt}
\begin{multicols*}{2}

\section*{I. LÝ THUYẾT CƠ BẢN}

\section*{Câu 1. Sự ra đời, phát triển và vai trò của triết học Mác -- Lênin}
\subsection*{Nội dung rút gọn}
\begin{itemize}
  \item Triết học Mác -- Lênin ra đời vào thế kỷ XIX do điều kiện lịch sử, xã hội và khoa học chín muồi.
  \item Kế thừa có chọn lọc: triết học cổ điển Đức, kinh tế chính trị học Anh, chủ nghĩa xã hội không tưởng Pháp, thành tựu khoa học tự nhiên.
  \item Mác, Ăngghen xây dựng nền tảng; Lênin phát triển thêm trong hoàn cảnh mới.
  \item Bước ngoặt cách mạng: không chỉ giải thích thế giới mà còn cải tạo thế giới bằng thực tiễn.
  \item Hiện nay, vẫn là thế giới quan giúp nhìn sự vật trong mối liên hệ, vận động và phát triển.
\end{itemize}
\textbf{Nhớ nhanh:} Ra đời từ thực tiễn + kế thừa tinh hoa + hướng tới cải tạo thế giới.

\section*{Câu 3. Mâu thuẫn lôgích và mâu thuẫn biện chứng}
\subsection*{Nội dung rút gọn}
\begin{itemize}
  \item \textbf{Mâu thuẫn lôgích}: lỗi trong tư duy; vừa khẳng định vừa phủ định cùng một điều.
  \item \textbf{Mâu thuẫn biện chứng}: mâu thuẫn có thật trong sự vật; các mặt đối lập vừa thống nhất vừa đấu tranh.
  \item Sự đấu tranh làm sự vật vận động, biến đổi và phát triển.
  \item Không có mâu thuẫn thì không có động lực phát triển.
  \item Trong AI: phát triển nhanh nhưng phải an toàn; hiệu quả lớn nhưng có rủi ro đạo đức; dữ liệu lớn nhưng bảo vệ quyền riêng tư; tự động hóa cao nhưng ảnh hưởng việc làm.
\end{itemize}
\textbf{Nhớ nhanh:} Mâu thuẫn lôgích = sai trong nghĩ; mâu thuẫn biện chứng = động lực phát triển.

\section*{Câu 6. Nhận thức, chân lý và tiến bộ}
\subsection*{Nội dung rút gọn}
\begin{itemize}
  \item \textbf{Nhận thức}: quá trình con người phản ánh thế giới khách quan; từ chưa biết đến biết, từ ít đến nhiều.
  \item \textbf{Chân lý}: tri thức đúng đắn, phù hợp hiện thực khách quan.
  \item Chân lý con người ở mỗi giai đoạn có tính tương đối, đúng trong điều kiện và phạm vi nhất định.
  \item Nhận thức là quá trình tích lũy ngày càng nhiều chân lý, vừa sửa sai vừa tiến bộ.
  \item Ví dụ: khoa học từng cho Mặt Trời quay quanh Trái Đất; y học sửa quan niệm về bệnh tật; công nghệ thay mô hình cũ bằng mô hình mới hiệu quả hơn.
\end{itemize}
\textbf{Nhớ nhanh:} Con người tiến bộ bằng tích lũy cái đúng + sửa cái sai.

\section*{Câu 7. Chân lý và thực tiễn (Lênin)}
\subsection*{Nội dung rút gọn}
\begin{itemize}
  \item Không thể chỉ giải quyết tranh luận lý thuyết; phải kiểm tra trong thực tiễn.
  \item \textbf{Thực tiễn là tiêu chuẩn của chân lý}.
  \item Thực tiễn bao gồm: sản xuất vật chất, hoạt động chính trị--xã hội, thực nghiệm khoa học.
  \item Lý thuyết đúng nếu: kết quả đúng trong thực tế, giúp giải quyết vấn đề, phù hợp hiện thực.
  \item Hiểu xong phải áp dụng, kiểm nghiệm và điều chỉnh.
\end{itemize}
\textbf{Nhớ nhanh:} Muốn biết đúng sai thật sự phải kiểm nghiệm trong thực tế.

\section*{II. LÝ THUYẾT NÂNG CAO}

\section*{Câu 2. Nguyên tắc phát triển}
\subsection*{Nội dung rút gọn}
\begin{itemize}
  \item Mọi sự vật đều vận động, biến đổi, không đứng yên.
  \item Phát triển: thay đổi từ thấp đến cao, từ đơn giản đến phức tạp, từ chưa hoàn thiện đến hoàn thiện hơn.
  \item Cơ sở: mâu thuẫn bên trong; sự đấu tranh giữa các mặt đối lập làm vật vận động.
  \item Phương pháp: nhìn sự vật trong vận động, phân biệt cái mới--cái cũ, ủng hộ cái mới tiến bộ.
  \item Mác vận dụng: xã hội không bất biến mà phát triển qua hình thái kinh tế--xã hội.
  \item AI: từ mô phỏng tư duy $\to$ hệ chuyên gia $\to$ học máy $\to$ học sâu $\to$ AI tạo sinh.
\end{itemize}
\textbf{Nhớ nhanh:} Muốn hiểu sự vật phải nhìn nó đang biến đổi như thế nào.

\section*{Câu 5. Ý thức phản ánh và sáng tạo}
\subsection*{Nội dung rút gọn}
\begin{itemize}
  \item Ý thức phản ánh: con người hiểu hiện thực khách quan.
  \item Ý thức sáng tạo: không chỉ phản ánh mà còn đặt mục tiêu, lựa chọn, dự đoán, tổ chức hành động.
  \item Con người biến ý tưởng (bản vẽ, kế hoạch) thành hiện thực qua thực tiễn.
  \item Thời đại tri thức, khoa học, công nghệ: vai trò ý thức càng lớn.
  \item AI là thành tựu tư duy con người; lợi ích: năng suất, học tập, y tế, nghiên cứu.
  \item Rủi ro: phụ thuộc công nghệ, thay đổi việc làm, lạm dụng, sai lệch dữ liệu, thiên vị thuật toán.
  \item Cần cân bằng: phát huy tích cực + kiểm soát rủi ro.
\end{itemize}
\textbf{Nhớ nhanh:} Ý thức không chỉ hiểu thế giới mà còn biến đổi thế giới.

\section*{Câu 8. Lý luận trở thành lực lượng vật chất}
\subsection*{Nội dung rút gọn}
\begin{itemize}
  \item Chỉ lời nói hay phê phán trên giấy chưa đủ thay đổi hiện thực.
  \item Cần lực lượng vật chất: hành động thực tế của con người.
  \item Lý luận đúng có sức mạnh rất lớn.
  \item Khi tư tưởng khoa học, tiến bộ thâm nhập quần chúng, được tiếp nhận và biến thành hành động, lý luận trở thành sức mạnh vật chất thật sự.
  \item Điều này cho thấy: vai trò lý luận đúng; vai trò quần chúng nhân dân; mối liên hệ tư tưởng--hành động.
  \item Cần: giáo dục, tổ chức và vận động quần chúng.
\end{itemize}
\textbf{Nhớ nhanh:} Lý luận đúng + quần chúng hành động = sức mạnh vật chất.

\section*{Câu 10. Các hình thái kinh tế--xã hội}
\subsection*{Nội dung rút gọn}
\begin{itemize}
  \item Xã hội không phát triển ngẫu nhiên mà theo quy luật khách quan.
  \item Các hình thái: cộng sản nguyên thủy, nô lệ, phong kiến, tư bản chủ nghĩa.
  \item Quá trình lịch sử--tự nhiên: diễn ra trong lịch sử thật nhưng tuân quy luật khách quan.
  \item Phương thức sản xuất giữ vai trò rất quan trọng.
  \item Khi lực lượng sản xuất phát triển, mâu thuẫn với quan hệ sản xuất cũ thúc xã hội sang hình thái mới.
  \item Học thuyết này là ``hòn đá tảng'' của duy vật lịch sử; giải thích xã hội từ cơ sở vật chất, phương thức sản xuất.
\end{itemize}
\textbf{Nhớ nhanh:} Xã hội phát triển theo quy luật khách quan, không ngẫu nhiên.

\section*{Câu 11. Mác tìm ra quy luật phát triển xã hội}
\subsection*{Nội dung rút gọn}
\begin{itemize}
  \item Darwin phát hiện quy luật phát triển thế giới sinh vật.
  \item Mác phát hiện quy luật phát triển xã hội loài người.
  \item Trước: lịch sử là chuỗi sự kiện rời rạc; giải thích bằng ý chí cá nhân.
  \item Mác: xã hội phát triển theo quy luật khách quan; phương thức sản xuất vật chất là nền tảng.
  \item Mối quan hệ: lực lượng sản xuất, quan hệ sản xuất, cơ sở hạ tầng, kiến trúc thượng tầng, đấu tranh giai cấp.
  \item Đóng góp: biến nghiên cứu xã hội từ cảm tính/suy đoán thành cách tiếp cận khoa học.
\end{itemize}
\textbf{Nhớ nhanh:} Darwin = sinh giới; Mác = lịch sử xã hội.

\section*{Câu 12. Tư duy lý luận và khoa học}
\subsection*{Nội dung rút gọn}
\begin{itemize}
  \item Muốn đứng đỉnh cao khoa học cần tư duy lý luận.
  \item Không thể chỉ kinh nghiệm rời rạc mà cần: khái quát, phân tích bản chất, phát hiện quy luật, liên hệ tri thức.
  \item Tư duy lý luận giúp: nhìn xa hơn, hiểu bản chất, dự báo xu hướng, xây dựng hệ thống tri thức.
  \item Mỗi lần khoa học đạt thành tựu mới, triết học phải thay đổi hình thức tồn tại.
  \item Phát minh khoa học (vật lý, sinh học, IT, AI) đặt ra câu hỏi mới cho triết học.
  \item Triết học chân chính: gắn bó khoa học, khái quát, tự đổi mới để không lạc hậu.
\end{itemize}
\textbf{Nhớ nhanh:} Khoa học phát triển cần tư duy lý luận; khoa học đổi thì triết học cũng đổi.

\columnbreak

\section*{III. ỨNG DỤNG THỰC TIỄN}

\section*{Câu 4. Cái mới, cái cũ và bài học CNTT}
\subsection*{Nội dung rút gọn}
\begin{itemize}
  \item \textbf{Cái mới}: phù hợp xu hướng phát triển, thay thế cái cũ, mở hướng tiến bộ.
  \item \textbf{Cái cũ}: từng phù hợp nhưng dần lạc hậu, cản trở phát triển.
  \item Đấu tranh: khó khăn, phức tạp, kéo dài vì cái cũ giữ vị trí, lợi ích quen thuộc; cái mới lúc đầu chưa hoàn thiện.
  \item Nếu cái mới phù hợp quy luật khách quan và nhu cầu phát triển, về lâu dài sẽ thắng.
  \item Trong CNTT: điện toán đám mây (từ bị nghi ngờ thành xu hướng lớn); AI (từ khó ứng dụng thành nhiều lĩnh vực); phương pháp phát triển mới (từ không tin tưởng được chấp nhận).
  \item Bài học: phát hiện cái mới tiến bộ, không bảo thủ, không nóng vội (cái mới cần thời gian), đánh giá bằng hiệu quả thực tế chứ không cảm tính.
\end{itemize}
\textbf{Nhớ nhanh:} Cái mới ban đầu yếu, nhưng nếu hợp xu thế cuối cùng vẫn thắng.

\section*{Câu 9. Thống nhất lý luận và thực tiễn}
\subsection*{Nội dung rút gọn}
\begin{itemize}
  \item \textbf{Lý luận}: khái quát kinh nghiệm thực tiễn thành tri thức có hệ thống.
  \item \textbf{Thực tiễn}: hoạt động vật chất có mục đích cải biến tự nhiên và xã hội.
  \item Thực tiễn là cơ sở, nguồn gốc của lý luận.
  \item Thực tiễn là mục đích và nơi kiểm tra tính đúng sai.
  \item Lý luận: soi đường, định hướng, giúp thực tiễn hiệu quả hơn.
  \item Yêu cầu: xuất phát thực tiễn, tôn trọng thực tiễn, học đi đôi hành, tránh lý thuyết suông, tránh kinh nghiệm mù quáng.
  \item Trong CNTT: học lập trình phải viết code; thiết kế phần mềm từ nhu cầu người dùng; sản phẩm kiểm thử--sửa lỗi--cải tiến; kiến thức lớp gắn dự án--bài tập--sản phẩm.
  \item Thống nhất lý luận--thực tiễn mới đạt hiệu quả cao.
\end{itemize}
\textbf{Nhớ nhanh:} Lý luận phải đi với làm; học phải gắn thực tế.

\end{multicols*}
\end{document}
